% Passages removed from paper.tex on 2026-10-21 (they were commented out).
% Kept here for reference; not \input by paper.tex.

% --------------------------------------------------------------------
% From Section 2.1 (sec:setup), after the Stokes kernel, eq. (stokes).
% Comparison with the closure of Foulger et al. (1992).
% --------------------------------------------------------------------
% \citet{foulger1992post} start from the same force balance but close it
% differently. Rather than coupling the elastic layer to an infinite viscous
% half-space through the non-local Stokes kernel (\ref{eq:stokes}), they place
% a viscous channel of finite thickness $H_v$ over a rigid base. Its basal
% traction is then treated as local, simple-shear drag,
% $\sigma_{xz}^{\mathrm{BOT}} \sim (\eta/H_v)\,\partial u/\partial t$, in place
% of the convolution in (\ref{eq:stokes}). Substituted into (\ref{eq:equil}),
% this closure collapses the governing equation to an ordinary, Fickian
% diffusion equation for the surface displacement, $\partial u/\partial t =
% \kappa\,\partial^2 u/\partial x^2$. The diffusivity $\kappa \propto E H
% H_v/\eta$ this introduces aggregates the elastic-layer properties with the
% assumed viscous-channel thickness $H_v$, rather than constraining a
% half-space viscosity directly. Reproducing a given deformation pattern
% therefore requires an a priori choice of $H_v$ at each site -- the
% limitation already noted in Section~\ref{sec:intro}. Section~\ref{sec:fourier}
% derives what the non-local closure (\ref{eq:stokes}) yields in its place: an
% infinite half-space, rather than a finite channel, replaces the ordinary
% Laplacian of Foulger et al.'s diffusion equation with the half-order
% fractional Laplacian $(-\Delta)^{1/2}$ of Equation~\ref{eq:fractional}. The
% two models therefore differ in the order of the governing diffusion
% equation itself, not only in its diffusivity.

% --------------------------------------------------------------------
% From Section 2.2 (sec:fourier), after eq. (fractional).
% Harmonic-extension (Caffarelli--Silvestre) origin of the half-order operator.
% --------------------------------------------------------------------
% This order is not a modelling choice. The fractional Laplacian admits a
% representation as a Dirichlet-to-Neumann map for a harmonic extension
% problem into an auxiliary half-space, weighted by $y^{1-2s}$; the weight
% vanishes, and the extension becomes the ordinary Laplace equation, only at
% $s=1/2$ \citep{caffarelli2007extension,daoud2022fractional}. That is the
% same half-space potential problem generating the Stokes kernel itself, so
% the half-order non-locality in (\ref{eq:fractional}) is forced by the
% geometry of the elastic-plate--viscous-half-space coupling, not fitted to
% the data. Solving the full two-dimensional Stokes flow beneath the plate,
% governed by the biharmonic rather than the harmonic equation, still
% recovers the same interface operator $-2\eta(-\Delta)^{1/2}$ exactly once
% the interface stress is extracted.

% --------------------------------------------------------------------
% From Section 2.2 (sec:fourier), after the previous block.
% --------------------------------------------------------------------
% This contrasts with classical Fickian diffusion: the half-order operator, a
% direct consequence of the non-local Stokes kernel, produces Lorentzian rather
% than Gaussian spatial profiles, with tails that decay as $1/x$ rather than
% exponentially.
